\documentclass[11pt]{article}
\usepackage{coursenotes}
\begin{document}
\notesheader{2}{Experiments}{Uncertainty from the design, decisions under it, and designing the next test}

\begin{decision}
An A/B test has run. The business needs to know whether the treatment pays for the population it will be
rolled out to, which is a different question from whether the effect is zero. Answering it takes an interval
on the right quantity, a test against the break-even, and a clear account of what randomisation does and
does not guarantee.
\end{decision}

\begin{goals}
  \item test the sharp null by randomisation and compute Neyman's standard error, and say why the usual SE
  is conservative;
  \item put an interval on net value and test against the break-even rather than zero;
  \item reduce variance with a pre-treatment covariate (CUPED) and size an experiment for a decision;
  \item choose the unit of randomisation, compute a design effect, and cluster standard errors correctly;
  \item recognise the threats that no formula repairs: post-assignment exclusions, sample-ratio mismatch,
  peeking, many metrics, interference;
  \item separate the effect of an offer (the ITT) from the effect of taking it up (the LATE), and say for whom the
  LATE holds.
\end{goals}

%=====================================================================================
\sectionone

\subsection{Setup and the estimand}

A completely randomised experiment draws $n_1$ of $n$ units for treatment and leaves $n_0 = n - n_1$ as
controls; $\hatt = \bar Y_1 - \bar Y_0$. In the \emph{design-based} frame the units and their potential
outcomes are fixed and the only randomness is the draw. The estimand is
$\tau_S = \frac1n\sum_i \{Y_i(1) - Y_i(0)\}$. Write
$S_d^2 = \frac{1}{n-1}\sum_i (Y_i(d) - \bar Y(d))^2$ for the spread of $Y(d)$ across units and $S_\tau^2$ for
the spread of $\tau_i$. In the \emph{superpopulation} frame units are an i.i.d.\ sample with
$\sigma_d^2 = \Var(Y(d))$, and the estimand is the population ATE.

\subsection{Identification after the draw}

Meeting~1 showed $\E_D[\hatt] = \tau_S$. Two conditions keep that true after the coin has been tossed.

\paragraph{Analyse as drawn.} Compare the arms the coin produced. Dropping units on anything measured
after assignment (``opened the app'', ``redeemed'', ``stayed to the end'') compares different kinds of unit
whenever treatment affects that variable. The difference in means over all assigned units identifies the
effect of \emph{assignment}, the intention-to-treat (ITT) effect, which is what a decision to send an offer
needs. When uptake is partial, the first stage $\E[D \mid Z=1] - \E[D \mid Z=0]$ is also identified, and
their ratio is an effect among compliers under further assumptions (Section~\ref{sec:noncompliance}).

\paragraph{The unit of assignment.} If unit $i$'s outcome depends on unit $j$'s treatment, the difference
in means compares two mixtures of spillover conditions, not ``everyone treated'' with ``no one treated''.
This is a failure of identification; no variance formula repairs it. The remedy is a unit of assignment
large enough to contain the spillover (Section~1.7 and Handout H4).

\paragraph{Sample-ratio mismatch.} Check the realised arm sizes against the planned allocation with a
chi-squared test, $\chi^2 = \sum_d (n_d - \E n_d)^2/\E n_d$ with one degree of freedom. A tiny $p$-value (in
practice below about $0.001$) means the draw or the logging is broken, for example a bot filter that removed
more users in one arm. Then the arms are no longer a random split, and no outcome analysis should proceed
until the cause is found.

\subsection{Uncertainty from the design}

\paragraph{Fisher: could the draw alone explain the result?} Fisher's approach holds the units fixed and asks
whether the draw alone could have produced the observed gap.
\begin{prop}{Randomisation test of the sharp null}{fisher}
Under the sharp null $Y_i(1) = Y_i(0)$ for every $i$, the observed outcomes are both potential outcomes,
so the distribution of any statistic over the $\binom{n}{n_1}$ equally likely draws is known exactly. The
$p$-value is the share of draws giving a statistic at least as extreme as the observed one.
\end{prop}
\begin{proof}
Under the null the outcome vector does not depend on the assignment; each draw maps to a known statistic
and has probability $1/\binom{n}{n_1}$.
\end{proof}
The test needs no model and no sampling assumption, and it respects the design automatically (clusters and
strata are re-drawn as they were drawn). Its limitation is its null: ``no effect for anyone'', not ``zero
average effect'' or ``net value below zero''.

\paragraph{Neyman: how far does one draw land?} Neyman's approach targets the average effect $\tau_S$ and asks
how far one draw's estimate can land from it.
\begin{prop}{Neyman's variance}{neyman}
Holding the potential outcomes fixed,
\[
  \Var_D(\hatt) = \frac{S_1^2}{n_1} + \frac{S_0^2}{n_0} - \frac{S_\tau^2}{n}.
\]
\end{prop}
\begin{proof}
Idea: each arm mean is a sample drawn without replacement, and the two arms partition the units, so their
means are negatively correlated. Without replacement, $\Var_D(\bar Y_d) = S_d^2/n_d - S_d^2/n$, and
$\Cov_D(\bar Y_1, \bar Y_0) = -S_{10}/n$ with $S_{10}$ the covariance of $Y(1)$ and $Y(0)$ across units.
Combining, the terms over $n$ sum to $(S_1^2 + S_0^2 - 2S_{10})/n = S_\tau^2/n$.
\end{proof}

\begin{prop}{The usual variance estimator is conservative}{conservative}
With $s_d^2$ the sample variance in arm $d$, $\widehat\Var = s_1^2/n_1 + s_0^2/n_0$ satisfies
$\E_D[\widehat\Var] = S_1^2/n_1 + S_0^2/n_0 \ge \Var_D(\hatt)$, with equality if and only if every unit has
the same effect.
\end{prop}
\begin{proof}
Within an arm the units are a simple random sample of the $n$ values of $Y(d)$, so $s_d^2$ is unbiased for
$S_d^2$. The gap is $S_\tau^2/n \ge 0$.
\end{proof}

The dropped term involves the joint distribution of $Y(1)$ and $Y(0)$, which no experiment reveals:
conservatism is the price of not knowing individual effects. Example: 1{,}000 customers made of 300 sure
things, 200 persuadables, 100 do-not-disturbs and 400 lost causes, split 500/500, have $\tau_S = 0.10$,
$S_1^2 = 250/999$, $S_0^2 = 240/999$ and $S_\tau^2 = 290/999$. The true standard deviation over draws is
$0.026$; the usual SE is $0.031$ (Figure~1).

\begin{prop}{Superpopulation variance}{superpop}
If units are i.i.d.\ draws and assignment is independent of potential outcomes,
$\Var(\hatt) = \sigma_1^2/n_1 + \sigma_0^2/n_0$, estimated by the same $\widehat\Var$.
\end{prop}
\begin{proof}
Given the assignment, $\bar Y_1$ and $\bar Y_0$ are means of independent draws from disjoint units, so their
variances add.
\end{proof}

The same formula has two readings: a conservative estimate of the design variance for the units in the test, and
the exact variance for a population they were sampled from. The second needs a sampling assumption; the
first needs only the draw log. The same SE is what OLS on a dummy reports when heteroskedasticity-robust
(HC2) standard errors are requested; the default homoskedastic SE is wrong when arms differ in size and
variance.

\begin{center}
\begin{tikzpicture}
\begin{axis}[width=0.78\linewidth, height=5.2cm, domain=0:0.2, samples=120, xlabel={$\hat\tau$ over draws},
  ylabel={density}, ytick=\empty, axis lines=left, legend style={draw=none, font=\footnotesize, at={(0.98,0.95)}, anchor=north east},
  xtick={0,0.05,0.10,0.15,0.20}, xticklabel style={/pgf/number format/fixed}]
\addplot[thick, sbgreen] {exp(-(x-0.10)^2/(2*0.026^2))/(0.026*sqrt(2*pi))};
\addplot[thick, dashed, noteblue] {exp(-(x-0.10)^2/(2*0.031^2))/(0.031*sqrt(2*pi))};
\legend{true randomisation distribution (SD $0.026$), normal with the usual SE ($0.031$)}
\end{axis}
\end{tikzpicture}

\small Figure 1. The usual SE describes a wider distribution than the one the draw actually produces.
\end{center}

\paragraph{Tests and intervals.} By a central limit theorem over draws, $(\hatt - \tau)/\widehat\SE$ is
approximately standard normal (conservatively so in the design-based frame). The interval
$\hatt \pm 1.96\,\widehat\SE$ covers the estimand in at least about 95\% of draws. It is the set of values the
test does not reject; zero is one point in it, with nothing special about it unless the decision is about
zero.

\subsection{Testing the decision rather than the effect}

Let a treated unit yield margin $m$ per purchase and cost $c$ per purchase (a coupon paid on redemption).
The net value per unit treated is $V = m\tau - c\,\mu_1 = (m-c)\mu_1 - m\mu_0$ (Meeting~1).

\begin{prop}{Inference on net value}{net}
$\hat V = (m-c)\bar Y_1 - m\bar Y_0$ is unbiased for $V$ over the draw, with (conservative) variance
estimate $(m-c)^2 s_1^2/n_1 + m^2 s_0^2/n_0$.
\end{prop}
\begin{proof}
$\hat V$ is a difference in means of the outcomes $(m-c)Y$ among treated and $mY$ among controls, so
Results~\ref{prop:conservative} and \ref{prop:superpop} apply to those outcomes.
\end{proof}

Testing $\hatt$ against a plug-in break-even $c\hat p_1/m$ treats an estimate as known; writing the decision
variable as a linear function of the arm means avoids this. When the cost is per unit treated, the net value
is $m\tau - c$ and the one-sided test is $H_0: \tau \le c/m$ with $z = (\hatt - c/m)/\widehat\SE$. In general,
\textbf{put the interval on the quantity the decision is made on.}

\subsection{Transport to another population}

An experiment on population $S$ identifies $\tau_S(x)$ for each $x$ in its support. A decision about a
population $T$ needs $\tau_T$, a different estimand.

\begin{assum}{Conditional effect invariance}{transport}
$\tau_T(x) = \tau_S(x)$ for every $x$ in $T$'s support, and that support is contained in $S$'s.
\end{assum}
\begin{prop}{Transport by reweighting}{transport}
Under Assumption~\ref{asm:transport}, $\tau_T = \sum_x \Pr_T(X=x)\,\tau_S(x)$.
\end{prop}
\begin{proof}
Iterated expectations in $T$, then the assumption; $\Pr_S(x) > 0$ wherever $\Pr_T(x) > 0$ by the support
condition, so every term is estimable.
\end{proof}
The assumption fails when the effect depends on something not in $X$ that differs between populations: the
season, the channel, the price level. A transported effect must also be held against the \emph{target's}
break-even, which moves with its composition. Types absent from the experiment cannot be reweighted into it.

\subsection{Designing the next experiment}

\paragraph{Fix it in writing first.} Unit of assignment and allocation; the outcome and its window,
measured for every assigned unit; the decision threshold; the sample size and stopping rule; logged exposure
and guardrail metrics. Then check the arm sizes on day one.

\begin{prop}{Sample size}{n}
With $n$ units per arm, a level-$\alpha$ two-sided test has approximate power $1-\beta$ against an effect
$\delta$ away from the null value when
\[
  n = \frac{(z_{1-\alpha/2} + z_{1-\beta})^2(\sigma_1^2 + \sigma_0^2)}{\delta^2}, \quad\text{rounded up.}
\]
For a one-sided test replace $z_{1-\alpha/2}$ by $z_{1-\alpha}$.
\end{prop}
\begin{proof}
With $\SE = \sqrt{(\sigma_1^2+\sigma_0^2)/n}$, the test rejects when $\hatt$ exceeds the null by
$z_{1-\alpha/2}\SE$; under the alternative this has probability $\Pr(Z > z_{1-\alpha/2} - \delta/\SE)$.
Setting it to $1-\beta$ gives $\delta = (z_{1-\alpha/2} + z_{1-\beta})\SE$; solve for $n$.
\end{proof}

Inverted, $\text{MDE} = (z_{1-\alpha/2} + z_{1-\beta})\SE$, about $2.8$ SEs at $\alpha = 0.05$ and 80\% power.
For a decision the null value is the break-even, so $\delta$ is the margin above break-even you would regret
missing; a test powered for ``is it zero?'' can be badly underpowered for ``does it pay?''. Halving the MDE
needs four times the sample.

\paragraph{Variance reduction with a pre-treatment covariate.} Let $X$ be fixed before assignment (last
month's spend, say) and define $Y^{\text{cv}} = Y - \theta(X - \bar X)$.

\begin{prop}{CUPED}{cuped}
For any fixed $\theta$, $\hatt^{\text{cv}} = \hatt - \theta(\bar X_1 - \bar X_0)$ is unbiased over the draw.
Within an arm, $\Var(Y - \theta X)$ is minimised at $\theta^\ast = \Cov(Y,X)/\Var(X)$, where it equals
$\Var(Y)(1-\rho^2)$, $\rho = \mathrm{Corr}(Y, X)$.
\end{prop}
\begin{proof}
$X$ is fixed before a random assignment, so $\E_D[\bar X_1 - \bar X_0] = 0$. The variance
$\Var(Y) - 2\theta\Cov(Y,X) + \theta^2\Var(X)$ is a quadratic minimised at $\theta^\ast$; substitute.
\end{proof}

In any one draw $\bar X_1 \ne \bar X_0$, so the adjusted estimate moves; that shift corrects chance imbalance.
Estimating $\theta$ from the pooled data costs nothing asymptotically. CUPED is regression adjustment; adding
the interaction $D \times (X - \bar X)$ (Lin's estimator) guarantees adjustment never hurts in large samples.
The required sample size scales by $1 - \rho^2$. If $X$ were affected by treatment, the correction would
subtract part of the effect.

\subsection{Clustered assignment}

When groups (stores, delivery zones, companies) are randomised, the design was randomised $J$ times, not once
per unit, and outcomes within a group are correlated.

\begin{prop}{Design effect}{deff}
For a cluster of $m$ units with common variance $\sigma^2$ and pairwise correlation $\rho_c$,
$\Var(\text{cluster mean}) = \tfrac{\sigma^2}{m}[1 + (m-1)\rho_c]$: the variance is inflated by
$\text{DEFF} = 1 + (m-1)\rho_c$ relative to $m$ independent units.
\end{prop}
\begin{proof}
$\Var(\tfrac1m\sum_j Y_j) = \tfrac{1}{m^2}[m\sigma^2 + m(m-1)\rho_c\sigma^2]$.
\end{proof}

With $J$ clusters of size $m$ the effective sample is $Jm/\text{DEFF}$; $\rho_c = 0.01$ and $m = 400$ give
$\text{DEFF} = 4.99$. Either compare cluster means ($J$ observations) or cluster standard errors at the level
of assignment. With fewer than about 40 clusters use small-sample corrections or randomisation inference over
cluster assignments. Choose the smallest unit that contains the spillover: the household for shared codes,
the store for staff behaviour, catchment clusters for customers who use nearby stores, time blocks with
washouts for offers that carry over.

\subsection{Peeking and many metrics}

The 5\% error rate of a test holds for one look at a sample size fixed in advance. Stopping at the first
$p < 0.05$ over ten equally spaced looks raises the false-positive rate under no effect to about 19\%.
Either look once, or use a method built for monitoring (group-sequential boundaries or always-valid tests,
Section~3). Likewise twenty metrics at 5\% yield about one false positive when nothing works. Declare one
primary metric tied to the decision; treat others as guardrails or exploration; when several tests drive the
decision, control the family-wise error (Holm's step-down procedure: order the $p$-values and compare the
$k$-th smallest with $\alpha/(K - k + 1)$, stopping at the first failure).

\subsection{Noncompliance: the ITT and the LATE}\label{sec:noncompliance}

Often the coin assigns an \emph{offer} and units decide whether to take it up. Write $Z \in \{0,1\}$ for
assignment and $D \in \{0,1\}$ for receipt. Each unit has \emph{potential receipts} $D(0), D(1)$ and potential
outcomes $Y(d, z)$; we observe $D = D(Z)$ and $Y = Y(D, Z)$. Randomising $Z$ delivers independence,
$\{Y(d,z), D(z) : d, z \in \{0,1\}\} \indep Z$. Two more assumptions are needed before anything can be said
about receipt.

\begin{assum}{Exclusion}{excl-nc}
$Y(d, 0) = Y(d, 1) \equiv Y(d)$: assignment affects the outcome only through receipt.
\end{assum}
\begin{assum}{Monotonicity}{mono-nc}
$D(1) \ge D(0)$ for every unit: nobody takes the treatment because they were \emph{not} offered it.
\end{assum}

\begin{prop}{Two effects of assignment}{itt}
Under random assignment alone, $\text{ITT} \equiv \E[Y \mid Z=1] - \E[Y \mid Z=0] = \E[Y(D(1),1) - Y(D(0),0)]$ and
the first stage $\pi \equiv \E[D \mid Z=1] - \E[D \mid Z=0] = \E[D(1) - D(0)]$.
\end{prop}
\begin{proof}
Given $Z = z$, $Y = Y(D(z), z)$ and $D = D(z)$; independence lets the conditioning be dropped.
\end{proof}

The ITT is the effect of the offer through every channel, with take-up as it happened. It is the number an
offer decision needs, and it needs neither exclusion nor monotonicity. Those are needed only to attribute an
effect to receipt itself.

Under monotonicity each unit is an \emph{always-taker} ($D(0) = D(1) = 1$), a \emph{never-taker}
($D(0) = D(1) = 0$) or a \emph{complier} ($D(0) = 0$, $D(1) = 1$); \emph{defiers} ($D(0) = 1$, $D(1) = 0$) are
ruled out. These are Meeting~1's response types, written for receipt instead of the outcome, and as there, no
unit's type is observed.

\begin{prop}{Type shares}{shares}
Under random assignment and monotonicity, $\pi_A = \Pr(D=1 \mid Z=0)$, $\pi_N = \Pr(D=0 \mid Z=1)$ and
$\pi_C = \pi$.
\end{prop}
\begin{proof}
With no defiers, $D = 1$ under $Z = 0$ only for always-takers and $D = 0$ under $Z = 1$ only for never-takers;
independence makes each arm's shares the population's. The rest are compliers, and $\pi = (\pi_A + \pi_C) - \pi_A$.
\end{proof}

\begin{prop}{The LATE theorem (Imbens and Angrist)}{late}
Under random assignment, exclusion and monotonicity, with $\pi > 0$,
\[
  \frac{\text{ITT}}{\pi} = \frac{\E[Y \mid Z=1] - \E[Y \mid Z=0]}{\E[D \mid Z=1] - \E[D \mid Z=0]}
  = \E[Y(1) - Y(0) \mid \text{complier}] \equiv \text{LATE}.
\]
\end{prop}
\begin{proof}
By exclusion the ITT is $\E[Y(D(1)) - Y(D(0))]$. For always- and never-takers $D(1) = D(0)$ and the difference
is zero; for compliers it is $Y(1) - Y(0)$; there are no defiers. So $\text{ITT} = \pi_C\,\text{LATE}$; divide by
$\pi = \pi_C > 0$.
\end{proof}

\begin{prop}{One-sided noncompliance}{onesided}
If no unit assigned to control can receive the treatment ($D(0) = 0$ for all), then under random assignment and
exclusion, $\text{ITT}/\pi = \E[Y(1) - Y(0) \mid D = 1]$: the ratio is the effect on those who received it.
\end{prop}
\begin{proof}
With $D(0) = 0$ there are no always-takers or defiers, so monotonicity holds, and the units with $D = 1$ are the
compliers assigned to treatment. By independence they have the same potential outcomes on average as all
compliers, so their average effect is the LATE.
\end{proof}

The ratio (the \emph{Wald estimator}) rescales the ITT and changes the population: from everyone assigned to the
compliers, a group the data choose. Two tempting alternatives fail. Comparing those who received the treatment
with those who did not (\emph{as treated}), or keeping only units that followed their assignment
(\emph{per protocol}), conditions on receipt, which assignment caused, so each compares different mixtures of
types (Section~1.2). Hold the ITT against a cost per offer and the LATE against a cost per receiver; when the
cost is paid only on receipt and noncompliance is one-sided, the two comparisons agree (Section~2, Step~6).
Instruments that are not randomised offers, and what failures of exclusion or monotonicity do to the ratio, are
in Handout H1.

\begin{pitfall}[what no standard error can fix]
Excluding units on post-assignment behaviour; reading a test with a sample-ratio mismatch; analysing a
store-randomised test with customer-level SEs; stopping when the dashboard turns green; reporting whichever
of twenty metrics moved.
\end{pitfall}

\checkpoint{A checkout test reports a 0.8-point lift with $p = 0.003$, using 30 metrics, peeked at daily, with
users randomised but the page cached by region. List every reason to doubt the decision, in order of how much
it matters.}

%=====================================================================================
\sectiontwo

\paragraph{Setting.} QuickBite, a food-delivery platform, tests a ¥10 discount on a customer's second order.
The outcome $Y$ is the customer's gross margin over the next 30 days, before the discount cost. The discount
costs ¥1.50 per customer offered on average (most do not redeem), so the break-even is $\tau^\ast = $ ¥1.50.
Customers are randomised 50/50, $n_1 = n_0 = 4{,}000$.

\paragraph{Step 1: checks.} Arm sizes match the plan. (In an earlier QuickBite test that planned a 50/50 split
of 20{,}000 and logged 10{,}240 and 9{,}760, $\chi^2 = 2 \times 240^2/10{,}000 = 11.52$, $p = 0.0007$: that test
was not read until a logging fault was fixed.)

\paragraph{Step 2: the raw estimate.} $\bar Y_1 = 21.80$, $\bar Y_0 = 19.90$, $s_1 = s_0 = 18$:
\[
  \hatt = 1.90, \qquad \widehat\SE = 18\sqrt{2/4{,}000} = 0.402, \qquad \text{95\% CI } [1.11, 2.69].
\]
Against zero $z = 4.72$; against the break-even, $z = (1.90 - 1.50)/0.402 = 0.99$, one-sided $p = 0.16$. The
discount certainly does something; whether it pays is unclear.

\paragraph{Step 3: CUPED.} Pre-period 30-day margin has $\rho = 0.7$ with $Y$ and $\hat\theta = 0.75$; by chance
$\bar X_1 - \bar X_0 = 0.08$. So $\hatt^{\text{cv}} = 1.90 - 0.75 \times 0.08 = 1.84$ and
$\widehat\SE \approx 0.402\sqrt{1 - 0.49} = 0.287$. The interval is $[1.28, 2.40]$ and $z = (1.84 - 1.50)/0.287 = 1.18$,
$p = 0.12$: narrower by 29\%, still not clear of the break-even.

\paragraph{Step 4: the follow-up's size.} To detect an effect ¥0.50 above break-even, one-sided at 5\% with
80\% power and CUPED variance $18^2 \times 0.51 = 165.2$ per arm:
$n = (1.645 + 0.842)^2 \times 2 \times 165.2/0.50^2 = 8{,}174.3$, so $8{,}175$ per arm. The current test's one-sided MDE above
break-even is $2.486 \times 0.287 = $ ¥0.71.

\paragraph{Step 5: had it been randomised by delivery zone.} With 500 customers per zone and $\rho_c = 0.01$,
$\text{DEFF} = 1 + 499 \times 0.01 = 5.99$: twenty zones (10{,}000 customers) carry the information of about
$1{,}670$ independent customers, and the SE must be clustered by zone.

\paragraph{Step 6: assignment versus receipt.} Only offered customers could use the discount, and 15\% of them
did: one-sided noncompliance with $\pi = 0.15$ (the cost of ¥1.50 per offer is $0.15 \times$ ¥10). The ITT of
¥1.84 is the effect of \emph{offering}, the number the offer decision needs. By
Result~\ref{prop:onesided}, the ratio $1.84/0.15 = $ ¥12.27 is the effect on customers who used the discount.
Held against the ¥10 each of them cost, it leaves ¥2.27, which is $(1.84 - 1.50)/0.15$: the same test as the ITT
against ¥1.50 per offer, rescaled by $1/0.15$. The ratio adds an interpretation (a customer who uses the
discount gains about ¥12 of margin) but not a new verdict, and its interval is $1/0.15$ times as wide.

\paragraph{Recommendation.} Do not roll out yet: the effect is real but not shown to exceed break-even. Run the
follow-up at about 8{,}200 per arm with CUPED and 30-day margin declared as the primary metric.

\begin{exercise}[computation]
QuickBite plans the follow-up by zone instead: zones of 400 customers, $\rho_c = 0.005$, CUPED variance as
above. How many zones per arm give the same power for a ¥0.50 margin above break-even?
\end{exercise}
\answer{$\text{DEFF} = 1 + 399 \times 0.005 = 2.995$. Customers needed per arm: $8{,}175 \times 2.995 \approx 24{,}480$, so
$24{,}480/400 \approx 61.2$, i.e.\ 62 zones per arm. (With so many zones the cluster-robust SE is reliable.)}

\begin{exercise}[judgement]
Marketing proposes to report the effect only among customers who opened the discount notification, ``because
the others never saw it''. Explain what this compares and when, if ever, it would equal an effect of the
discount.
\end{exercise}
\answer{Opening is measured after assignment and the notification itself changes who opens, so treated
openers are compared with control customers who opened something different (or are a different mix of
people). The comparison carries selection bias. It equals an effect only if opening were unaffected by
assignment and unrelated to potential outcomes, which is implausible. Report the ITT; for the effect among
those who engage, divide the ITT by the share who engage when offered (Result~\ref{prop:late}).}

%=====================================================================================
\sectionthree

\begin{extension}[More on randomisation inference]
Fisher tests extend to any statistic (rank sums, trimmed means) and to sharp nulls of a constant effect
$\tau_0$; inverting them over $\tau_0$ gives a confidence set. Under complete randomisation Neyman's
variance can be tightened using bounds on $S_\tau^2$ (Aronow, Green and Lee, 2014).
\end{extension}

\begin{extension}[Stratified designs and ratio metrics]
Randomising within strata of a strong predictor removes chance imbalance by design; the estimator is the
stratum-weighted difference and its variance sums the within-stratum variances. Ratio metrics (revenue per
session) have numerator and denominator in both arms; their standard error needs the delta method, and
randomising users while analysing sessions needs clustering by user.
\end{extension}

\begin{extension}[Beyond randomised offers]
When the offer is not randomised but something else shifts treatment (a rule, a lottery, an outage), the same
ratio identifies a LATE only if that shifter is as good as random, which must now be argued. Handout H1 covers such
\emph{instruments}: two-stage least squares with covariates, weak instruments, what failures of exclusion and
monotonicity do, and how to describe the compliers.
\end{extension}

\begin{extension}[Sequential testing]
Group-sequential designs spend the error rate over planned looks (O'Brien--Fleming, Pocock). The mixture
sequential probability ratio test gives always-valid $p$-values that may be checked at any time.
\end{extension}

\subsection*{Annotated reading}
\begin{readinglist}
\reading{Imbens and Rubin (2015), \emph{Causal Inference for Statistics, Social, and Biomedical
Sciences}}{Fisher and Neyman inference in the design-based frame, with proofs; noncompliance in experiments}{Ch.\ 5--6, 9, 23--24}{2}
\reading{Angrist, Imbens and Rubin (1996), ``Identification of causal effects using instrumental variables'',
\emph{JASA}}{Compliance types and the LATE in potential outcomes}{Sections 1--4}{2}
\reading{Kohavi, Tang and Xu (2020), \emph{Trustworthy Online Controlled Experiments}}{The practitioner's
checklist: SRM, guardrails, peeking, ratio metrics}{Ch.\ 3, 17, 21}{1}
\reading{Deng, Xu, Kohavi and Walker (2013), ``Improving the sensitivity of online controlled experiments'',
\emph{WSDM}}{The original CUPED paper}{All}{1}
\reading{Lin (2013), ``Agnostic notes on regression adjustments to experimental data'', \emph{Annals of
Applied Statistics}}{Why the interacted adjustment never hurts asymptotically}{Sections 1--3}{3}
\reading{Abadie, Athey, Imbens and Wooldridge (2023), ``When should you adjust standard errors for
clustering?'', \emph{QJE}}{Cluster at the level of assignment, and why}{Sections 1--3}{2}
\reading{Johari, Koomen, Pekelis and Walsh (2022), ``Always valid inference'', \emph{Operations
Research}}{Continuous monitoring without inflating errors}{Sections 1--4}{3}
\reading{Handout H4, \emph{Interference and marketplace experiments}}{What to do when units affect each
other}{All}{2}
\end{readinglist}

\end{document}
